% sections/01-reconocimiento.tex
\section{Fase de reconocimiento}

\subsection{Selección del objetivo}

Para el desarrollo de esta primera fase se seleccionó el dominio \texttt{urosario.edu.co}, sobre el cual se ejecutaron las herramientas en línea dispuestas por la guía, dejando para el final las consultas DNS avanzadas, que la guía pide explícitamente sobre un dominio diferente.

\porque{dominio elegido}{%
  Cambia \texttt{urosario.edu.co} por el dominio que hayan escogido realmente
  y ajusta las descripciones de las capturas en consecuencia: lo importante es
  que sea un dominio público, que el mismo aparezca en ping, traceroute y
  Whois, y que el de la subsección de DNS (MxToolbox) sea uno diferente.%
}

\subsection{Ping en línea}

Primeramente se accedió desde el navegador a la herramienta de red en línea \cite{networktools}, se ingresó a la opción de ping y, con el dominio seleccionado en el cuadro de texto, se inició la prueba:

\captura{ping.png}{Resultado del ping en línea sobre el dominio seleccionado}

\pregunta{¿Para qué sirve el comando Ping y qué ventajas cree que se tienen al realizar un ping desde Internet?}
\begin{respuesta}
\end{respuesta}

\porque{base para la respuesta}{%
  Sirve para verificar alcanzabilidad y latencia: envía paquetes ICMP echo
  request y espera los echo reply, midiendo el tiempo de ida y vuelta. La
  ventaja de hacerlo desde Internet (una herramienta externa) es que la prueba
  sale desde fuera de la red interna, de forma que confirma que el dominio
  resuelve y es alcanzable públicamente (algo relevante, por ejemplo, para
  saber si un cortafuegos bloquea el ICMP), además de dar una vista externa
  de la ruta y la latencia sin depender de la propia conexión. Mencionar
  también la limitación: muchos servidores bloquean ICMP por seguridad.%
}

\subsection{Traceroute}

Seguidamente se seleccionó la opción traceroute de la misma herramienta y, con el dominio seleccionado, se inició el trazado de la ruta:

\captura{traceroute.png}{Resultado del traceroute sobre el dominio seleccionado}

\pregunta{¿Para qué sirve el comando Tracert?}
\begin{respuesta}
\end{respuesta}

\porque{base para la respuesta}{%
  Muestra la ruta (saltos intermedios, routers) que siguen los paquetes hasta
  el destino, con el tiempo de cada salto: se usa para diagnosticar dónde se
  degrada o se corta una conexión. Con la vista externa que da la herramienta
  en línea se evidencian además los puntos de intercambio y proveedores por
  los que entra el tráfico al país o a la red del objetivo, información de
  infraestructura útil en la fase de reconocimiento.%
}

\subsection{Búsqueda Whois}

Con la opción de búsqueda Whois de la misma herramienta y el dominio seleccionado se obtuvieron los datos de registro del dominio.

\pregunta{¿Para qué sirve el protocolo WhoIs?}
\begin{respuesta}
\end{respuesta}

\porque{base para la respuesta}{%
  Es el protocolo de consulta de las bases de datos de registro de dominios:
  entrega información como registrante, organización, fechas de creación y
  vencimiento, servidores de nombres y datos de contacto del administrador,
  información que en la fase de reconocimiento permite perfilar a la
  organización dueña del dominio (proveedores, rangos asociados, correos de
  contacto que alimentan, por ejemplo, campañas de phishing dirigido).%
}

\subsection{Consultas DNS avanzadas (MxToolbox)}

Para las consultas a nivel de DNS públicos se accedió a la herramienta MxToolbox \cite{mxtoolbox} y se seleccionó un dominio diferente al usado hasta el momento (\texttt{gmail.com}); la primera consulta, con el tipo por defecto (A), generó la siguiente respuesta:

\captura{dns-consulta.png}{Consulta DNS por defecto (tipo A) sobre el segundo dominio seleccionado}

\pregunta{¿Qué respuesta generó la consulta DNS al dominio que usted seleccionó?}
\begin{respuesta}
\end{respuesta}

\porque{base para la respuesta}{%
  Describe la salida: el registro A resuelve el dominio a una o varias
  direcciones IP (para gmail.com se esperan las IPs de Google), mostrando
  además el TTL de cada registro; esa dirección IP es la que un resolutor
  entrega a cualquier cliente que quiera conectar con el dominio. Ajusta la
  descripción al dominio que hayan usado.%
}

Usando el mismo dominio se cambió el tipo de consulta por MX:

\captura{dns-mx.png}{Consulta DNS de tipo MX sobre el segundo dominio seleccionado}

\pregunta{¿Qué registros se muestran ahora? Adjunte la respuesta.}
\begin{respuesta}
\end{respuesta}

\porque{base para la respuesta}{%
  Los registros MX (mail exchanger) listan los servidores de correo
  encargados de recibir el correo del dominio, cada uno con su prioridad
  (para gmail.com se espera \texttt{gmail-smtp-in.l.google.com} y los
  servidores alternativos con prioridades 5, 10, 20\dots{}): de cara al
  reconocimiento, revelan el proveedor de correo y la infraestructura
  asociada.%
}

\pregunta{Seleccione dos tipos de consulta (query type) que no se hayan visto hasta el momento y explíquelos brevemente.}
\begin{respuesta}
\end{respuesta}

\porque{sugerencia de tipos}{%
  Dos que rinden bien en MxToolbox: \textbf{TXT:} registros de texto
  arbitrarios, usados hoy sobre todo para las políticas de correo (SPF,
  DKIM, DMARC) y verificaciones de dominio; \textbf{SOA:} el registro de
  autoridad de la zona, con el servidor DNS maestro, el correo del
  administrador y los temporizadores de replicación. Otras opciones: NS
  (servidores de nombres de la zona), CNAME (alias de un nombre hacia otro),
  PTR (resolución inversa, de IP a nombre), AAAA (direcciones IPv6).%
}

\subsection{Listas negras (BLACKLIST)}

Finalmente, desde la página principal de MxToolbox se seleccionó la opción BLACKLIST y se consultó el mismo dominio:

\captura{blacklist.png}{Resultado de la búsqueda del segundo dominio en listas negras}

\pregunta{¿Qué son las listas negras? Dé ejemplos de otros sitios que puedan ser usados para hacer búsquedas de listas negras.}
\begin{respuesta}
\end{respuesta}

\porque{base para la respuesta}{%
  Son bases de datos públicas (DNSBL, \eng{blocklists}) donde se registran
  direcciones IP o dominios asociados a spam, malware o actividad maliciosa,
  consultadas por servidores de correo y firewalls para rechazar tráfico.
  Ejemplos de sitios para consultarlas: Spamhaus \cite{spamhaus}, Barracuda
  \cite{barracuda}, SORBS, SpamCop y la propia MXToolbox; mencionar qué
  significa para una organización terminar listada (correo rechazado) y que
  la salida de la consulta muestra en cuántas listas aparece el dominio
  (cero, si tiene buena reputación).%
}
